\documentclass[11pt,a4paper]{article}
\usepackage{amsmath,amssymb}
\usepackage{listings}
\usepackage{hyperref}
\usepackage[margin=1in]{geometry}
\usepackage{fancyhdr}

\pagestyle{fancy}
\fancyhf{}
\fancyhead[L]{Module 02: Variables and Data Types}
\fancyhead[R]{C Programming Course}
\fancyfoot[C]{\thepage}

\lstset{
  language=C,
  basicstyle=\ttfamily\small,
  keywordstyle=\bfseries,
  commentstyle=\itshape,
  numbers=left,
  numberstyle=\tiny,
  frame=single,
  breaklines=true,
  tabsize=4
}

\title{Module 02: Variables and Data Types}
\author{C Programming Course}
\date{}

\begin{document}
\maketitle
\thispagestyle{fancy}

\section{Variables and Constants}
A variable is a named memory location. Constants are declared with \texttt{const}
or \texttt{\#define}.

\begin{lstlisting}
#include <stdio.h>

int main(void) {
    int count = 10;
    const double PI = 3.14159;
    printf("count = %d, PI = %.5f\n", count, PI);
    return 0;
}
\end{lstlisting}

\section{Primitive Data Types}
\begin{center}
\begin{tabular}{lll}
\hline
\textbf{Type} & \textbf{Size (typical)} & \textbf{Format} \\
\hline
\texttt{char}   & 1 byte  & \texttt{\%c} \\
\texttt{int}    & 4 bytes & \texttt{\%d} \\
\texttt{float}  & 4 bytes & \texttt{\%f} \\
\texttt{double} & 8 bytes & \texttt{\%lf} \\
\hline
\end{tabular}
\end{center}

\section{Type Conversion and Casting}
\subsection{Implicit Conversion}
The compiler automatically converts smaller types to larger types in expressions.

\subsection{Explicit Casting}
\begin{lstlisting}
#include <stdio.h>

int main(void) {
    int a = 7, b = 2;
    double result = (double)a / b;
    printf("7 / 2 = %.2f\n", result);
    return 0;
}
\end{lstlisting}

\section{Operators}
\subsection{Arithmetic Operators}
\texttt{+}, \texttt{-}, \texttt{*}, \texttt{/}, \texttt{\%} (modulo).

\subsection{Relational Operators}
\texttt{==}, \texttt{!=}, \texttt{<}, \texttt{>}, \texttt{<=}, \texttt{>=}.

\subsection{Logical Operators}
\texttt{\&\&} (AND), \texttt{||} (OR), \texttt{!} (NOT).

\subsection{Bitwise Operators}
\texttt{\&}, \texttt{|}, \texttt{\^}, \texttt{\~}, \texttt{<<}, \texttt{>>}.

\begin{lstlisting}
#include <stdio.h>

int main(void) {
    unsigned int flags = 0x0F;
    flags = flags & 0x03;  /* mask lower 2 bits */
    printf("flags = 0x%02X\n", flags);
    return 0;
}
\end{lstlisting}

\section{Pre/Post Increment and Decrement}
\begin{lstlisting}
#include <stdio.h>

int main(void) {
    int x = 5;
    printf("%d\n", x++); /* prints 5, then x becomes 6 */
    printf("%d\n", ++x); /* x becomes 7, then prints 7 */
    return 0;
}
\end{lstlisting}

\end{document}
